\scanpage{449}
\mspage{361}
D'autre part, il est clair que $I$ est petite, je dis que $I$ est stable par $\underrightarrow{\lim}$ filtrantes de cardinal $\leq\pi$, et que le foncteur $\deleted{x}J'\longmapsto x_{J'}:I\unsure{\text{---}|\longrightarrow}M$ y commute, enfin que les $x_{J'}$ sont dans $M_\pi$. Il reste à voir que $I$ est grde devant $\pi$, i.e. que \deleted{pour les sous-catégories} si on a $(J'_{\unsure{r|i}})_{r\in A}$, \inserted{avec} $\mathrm{Card}\,A\leq\pi$, les $J'_r$ dans $I$, alors $\exists\,J'\in I$ qui majore les $J'_r$. Soit \deleted{$K$} \inserted{$K$} \deleted{soit} la sous-catégorie \deleted{str. pleine} de $J$ engendrée par les $J'_r$, on voit \deleted{tout} \deleted{\unsure{facilement}} \deleted{qui est telle que $\mathrm{Ob}\,K=\bigcup\mathrm{Ob}\,J'_r$, \unsure{tout}} que $\mathrm{Card}\,K\leq\pi$. \unsure{restant|revenant} à prouver :

% slanted
\margin{\uline{NB} Pour la démonstration de 4.13.37, il suffit de prouver 4.13.38 dans le cas où $J$ est ordonné, auquel cas les \unsure{contorsions|complications|considérations} de la fin sont inutiles.}
\uline{Lemme} 4.13.38. Soient $J$ une catégorie gr\unsure{de|ande} devant $\pi$, $K$ une sous-catégorie de $J$ \deleted{avec $\mathrm{Card}$} de cardinal $\leq\pi$.

Alors il existe une sous-catégorie $K'$ \deleted{\unsure{cont}} \uline{filtrante} de $J$, contenant $K$ de cardinal $\leq\pi$, [\deleted{Si $K$ contient tous les objets \unsure{maxim$^{\rm x}$} de $J$,} \deleted{On peut $=$ préciser que} \deleted{$K'=K\cup\{e\}$ tel que $K'$ ait un objet final dans $K$} (Si on suppose que l'on \unsure{n'est|n'ait} pas, \deleted{on peut choisir \unsure{$K'$} tout entier \unsure{égal}} ds le cas où $J$ contient un objet maximal, et que $K$ contienne tous les objets maximaux de $J$],
% archive sheet number: 448
